问题二更像是在搭建一个“珍稀物种响应模块”。题面附件中的事实需要逐项映射到参数：2022 年公报对应种群初值和监测基准；附件 1 中的航道阻隔、鱼道连通与栖息地破碎化对应 \(\delta_B\)；附件 2 中的放流、补充和食源恢复对应 \(R_S\)；\(e_D,m_D,H,u\) 则分别反映繁殖效率、自然死亡、捕捞/扰动和污染压力。只有把这些事实置于同一框架，才能说明禁渔收益如何经由食物网和生境连通性被放大或削弱，并解释在相同背景下，江豚和长江鲟为何呈现不同的恢复曲线\cite{report2022}。

上述事实映射还意味着，问题二必须把“自然恢复”和“工程补偿”分开解释。江豚更多表现为食源改善后的数量响应；长江鲟除食源外，还受到洄游通道是否连续、放流是否提供额外补充以及污染是否抵消恢复收益的影响。本问的重点不在于比较两个物种绝对值的大小，而在于考察不同约束条件下政策收益保留和损失的程度。这层关系把讨论从“珍稀物种是否增加”推进到“珍稀物种响应机制如何分化”。

表格的作用也在于让读者一眼看出每个参数到底来自哪条事实链，而不是把变量当成纯数学符号使用。问题二的证据链也由此变得清楚：上游是禁渔和食源恢复，中间由通道、放流与污染加以修正，下游才是珍稀物种的数量变化。
